\chapter{The Riemann Hypothesis}\label{conjectures:ch:riemann}

\begin{center}\itshape
It is very probable that all roots are real. One would of course wish for a rigorous proof of this; I have, however, after some fleeting vain attempts,
provisionally put aside the search for one.
\upshape(Bernhard Riemann, 1859, in our translation)
\end{center}

In 1859 Bernhard Riemann, newly elected a corresponding member of the Berlin Academy, sent the Academy the customary report on his research. It was
eight pages long, the only paper he ever wrote on number theory, and it contained an exact formula for the number of primes up to $x$, a sketch of
the functional equation of $\zeta$, a count of its zeros, and the sentence of the epigraph. The ``roots'' are those of $\xi(\frac12+it)$ as a function of
$t$; ``real'' means that the zeros of $\zeta$ in the critical strip lie on the line $\operatorname{Re}s=\frac12$. That is the Riemann hypothesis.

Riemann did more than guess. In 1932 Carl Ludwig Siegel went through the notes Riemann left in the G\"ottingen library and found there a formula for
computing $\zeta$ on the critical line, far better than anything published since, and computations of the first zeros. The \emph{Riemann--Siegel
formula} is still at the heart of every large computation. Hilbert put the hypothesis on his list of problems in 1900; the Clay Mathematics Institute
put it on its list of Millennium Problems in 2000, with a prize of a million dollars. In 2021 David Platt and Timothy Trudgian verified that every zero
with $0<\operatorname{Im}\rho\le3\cdot10^{12}$ lies on the line.

Tome~5 proved what the hypothesis says about primes: by von Koch's theorem (Tome~5, Theorem~7.10) it is equivalent to
$\psi(x)=x+O(\sqrt x\log^2x)$. This chapter looks at the hypothesis from two other sides. First the computations: how one checks zeros, and a pattern
called Gram's law that holds at every Gram point up to $g_{125}$ and then fails. Then a surprising reformulation in terms of the sum of divisors: Gronwall's
theorem, which we prove, says that $\sigma(n)$ can be as large as $e^\gamma n\log\log n$, and Robin's theorem says that the Riemann hypothesis is
\emph{exactly} the statement that it never gets larger, once $n>5040$.

\begin{think}
The number $5040=7!$ has $60$ divisors. Is $\sigma(5040)$ larger or smaller than $5040\cdot\log\log5040\cdot e^\gamma\approx19\,237$? What kind of $n$
make $\frac{\sigma(n)}n$ large?
\end{think}

\section{Counting zeros on the line}

We recall the facts from Tomes~4 and~5. The completed function $\xi(s)=\frac12s(s-1)\pi^{-s/2}\Gamma(\frac s2)\zeta(s)$ is entire, satisfies
$\xi(s)=\xi(1-s)$, and its zeros are the non-trivial zeros $\rho=\beta+i\gamma$ of $\zeta$ (Tome~5, Proposition~6.2). On the critical line,
\[
\xi\Bigl(\frac12+it\Bigr)=-\frac12\Bigl(\frac14+t^2\Bigr)\pi^{-1/4}\Bigl|\Gamma\Bigl(\frac14+\frac{it}2\Bigr)\Bigr|Z(t),\qquad Z(t)=e^{i\vartheta(t)}\zeta\Bigl(\frac12+it\Bigr),
\]
where $\vartheta(t)=\operatorname{Im}\log\Gamma(\frac14+\frac{it}2)-\frac t2\log\pi$. \emph{Hardy's function} $Z$ is real-valued, with $|Z(t)|=|\zeta(\frac12+it)|$
(Tome~4, Chapter~10), so every sign change of $Z$ is a zero on the line. The number $N(T)$ of zeros with $0<\gamma\le T$ satisfies
\[
N(T)=\frac{\vartheta(T)}\pi+1+S(T),\qquad S(T)=O(\log T),\qquad\vartheta(T)=\frac T2\log\frac T{2\pi}-\frac T2-\frac\pi8+O\Bigl(\frac1T\Bigr)
\]
(Tome~5, Theorem~6.9). To verify the hypothesis up to height $T$, one counts the sign changes of $Z$ on $]0,T]$ and checks that there are $N(T)$ of
them. Turing's method, explained in Tome~5, turns this into a rigorous procedure.

\paragraph{Gram points.} For $t\ge7$ the function $\vartheta$ is increasing. The \emph{Gram point} $g_n$ is the solution of $\vartheta(g_n)=n\pi$, for
$n\ge-1$; so $g_{-1}\approx9.67$, $g_0\approx17.85$, $g_1\approx23.17$. At a Gram point $e^{i\vartheta}=(-1)^n$, so
\[
\zeta\Bigl(\frac12+ig_n\Bigr)=(-1)^nZ(g_n)\quad\text{is real}.
\]
In 1903 J\o rgen Gram observed that the zeros of $Z$ and the Gram points tend to alternate, which means that $(-1)^nZ(g_n)>0$, that is
$\operatorname{Re}\zeta(\frac12+ig_n)>0$. There is a reason: the Riemann--Siegel formula writes $\zeta(\frac12+it)$ as a short sum
$\sum_{k\le\sqrt{t/2\pi}}k^{-1/2-it}$ plus a similar sum and a small error, and the first term of the sum is $1$, so the real part has a head start. In 1925
J.~I.~Hutchinson named the pattern \emph{Gram's law} and found where it first fails: at $g_{126}\approx282.455$, and again at $g_{134}\approx295.584$
(Figure~\ref{conjectures:fig:gram}). Gram's law is true for a positive proportion of $n$, and false for a positive proportion too (Timothy Trudgian, 2011). It is a good example of a
pattern that is \emph{often} right for a structural reason, and the computer programs that verify the hypothesis use it as a first guess and repair the
exceptions.

\begin{figure}[ht]
\centering
\fig{conjectures-ch04-gram}
\caption{Hardy's function $Z(t)$ for $270\le t\le300$, with the Gram points $g_{119},\dots,g_{136}$ on the axis. Gram's law $(-1)^nZ(g_n)>0$ fails at
$g_{126}$ and $g_{134}$ (red): there $Z$ has the ``wrong'' sign, the interval before is empty of zeros, and a neighbouring interval contains two. There
are $18$ zeros in the window, all simple and on the line.}
\label{conjectures:fig:gram}
\end{figure}

\paragraph{The record.} Gram checked the first $15$ zeros in 1903; Hutchinson $138$ in 1925; Edward Titchmarsh $1041$ in 1935, with the Riemann--Siegel
formula and a team using tabulating machines; Alan Turing pushed to height $1540$ in 1953, and Derrick Lehmer to $25\,000$ zeros in 1956. With
electronic computers the counts exploded: $1.5\cdot10^9$ zeros by Jan van de Lune, Herman te Riele and Dik Winter in 1986, $10^{13}$ by Xavier Gourdon in
2004, and height $3\cdot10^{12}$ (about $1.2\cdot10^{13}$ zeros) by Platt and Trudgian in 2021. Andrew Odlyzko computed billions of zeros near the
$10^{22}$-nd, all on the line; Appendix~E looks at their spacings.

\begin{remark}[on the line]\label{conjectures:rem:hardy}
Hardy proved in 1914 that infinitely many zeros lie on the line. Atle Selberg proved in 1942 that a positive proportion do, Norman Levinson in 1974
that more than a third do, and Brian Conrey in 1989 that more than two fifths do. These results are proved in a later tome. Not one zero is known off the line,
and not one zero is known to be multiple.
\end{remark}

\section{Gronwall's theorem}

We now turn to the sum of divisors $\sigma(n)=\sum_{d\mid n}d$. For $n=\prod p^{a_p}$,
\[
\frac{\sigma(n)}n=\sum_{d\mid n}\frac1d=\prod_{p\mid n}\Bigl(1+\frac1p+\dots+\frac1{p^{a_p}}\Bigr)=\prod_{p\mid n}\frac{1-p^{-a_p-1}}{1-p^{-1}}.
\]
Each factor is less than $(1-\frac1p)^{-1}$, so $\frac{\sigma(n)}n<\frac n{\varphi(n)}$. How large can this ratio be? The answer is a theorem of Thomas
Gronwall, from 1913, and the proof is Mertens' third theorem.

\begin{theorem}[Gronwall, 1913]\label{conjectures:thm:gronwall}
For $n\ge3$ let $G(n)=\dfrac{\sigma(n)}{n\log\log n}$. Then
\[
\limsup_{n\to\infty}G(n)=e^\gamma=1.7810724\ldots
\]
\end{theorem}
\begin{proof}
We use Mertens' theorem $\prod_{p\le y}(1-\frac1p)^{-1}=e^\gamma\log y\,(1+O(\frac1{\log y}))$ (Tome~5, Theorem~1.6), and the prime number theorem in the
form $\theta(y)=\sum_{p\le y}\log p\sim y$ (Tome~5, Theorem~5.5).

\emph{Upper bound.} Let $n\ge3$, let $k=\omega(n)$ be its number of prime factors, and let $y$ be the $k$-th prime. The $i$-th smallest prime factor of $n$
is at least the $i$-th prime, and $(1-\frac1p)^{-1}$ decreases with $p$, so
\[
\frac{\sigma(n)}n<\prod_{p\mid n}\Bigl(1-\frac1p\Bigr)^{-1}\le\prod_{p\le y}\Bigl(1-\frac1p\Bigr)^{-1}=:M(y).
\]
Also $n\ge\prod_{p\le y}p=e^{\theta(y)}$, so $\log\log n\ge\log\theta(y)$. Fix $\varepsilon>0$. Since $M(y)\sim e^\gamma\log y$ and $\log\theta(y)=\log y+o(1)$,
there is $Y$ such that $M(y)\le(1+\varepsilon)e^\gamma\log\theta(y)$ for $y\ge Y$; for such $n$, $G(n)\le(1+\varepsilon)e^\gamma$. If $y<Y$, then
$\frac{\sigma(n)}n<M(Y)$ and $G(n)<\frac{M(Y)}{\log\log n}$, which is below $e^\gamma$ for $n$ large. So $G(n)\le(1+\varepsilon)e^\gamma$ for all large $n$.

\emph{Lower bound.} For $x\ge3$ put $n_x=\prod_{p\le x}p^{a_p}$, where $a_p\ge1$ is the largest exponent with $p^{a_p}\le x^2$. Then $p^{a_p+1}>x^2$ for every
$p\le x$, so by Bernoulli's inequality
\[
\frac{\sigma(n_x)}{n_x}=\prod_{p\le x}\frac{1-p^{-a_p-1}}{1-p^{-1}}\ge\Bigl(1-\frac1{x^2}\Bigr)^{\pi(x)}M(x)\ge\Bigl(1-\frac{\pi(x)}{x^2}\Bigr)e^\gamma\log x\,\Bigl(1+O\Bigl(\frac1{\log x}\Bigr)\Bigr).
\]
On the other hand $\log n_x=\sum_{p\le x}a_p\log p\le2\pi(x)\log x\le3x$ for large $x$, by the prime number theorem, so $\log\log n_x\le\log x+\log3$.
Hence $G(n_x)\ge e^\gamma(1+o(1))$ as $x\to\infty$, and $n_x\to\infty$.
\end{proof}

The extremal numbers in the proof are products of all the small primes with decreasing exponents, numbers like $5040=2^4\cdot3^2\cdot5\cdot7$. Before
Gronwall, Ramanujan had studied such numbers in his 1915 memoir on highly composite numbers. Part of the memoir was cut for lack of paper during the war,
and in it Ramanujan proved that \emph{if} the Riemann hypothesis holds, then $G(n)<e^\gamma$ for all large $n$; the missing part was published only in
1997.

\section{Colossally abundant numbers}

To find the $n$ with the largest $G(n)$, Leonidas Alaoglu and Paul Erd\H{o}s (1944) trade $\sigma(n)/n$ against the size of $n$ with a parameter.

\begin{proposition}[Alaoglu--Erd\H{o}s]\label{conjectures:prop:ca}
Fix $\varepsilon>0$ and let $f(n)=\dfrac{\sigma(n)}{n^{1+\varepsilon}}$.
\begin{enumerate}[label=(\alph*)]
\item For each prime $p$ and $a\ge0$, $f(p^{a+1})\ge f(p^a)$ if and only if $p^{1+\varepsilon}\le\dfrac{\sigma(p^{a+1})}{\sigma(p^a)}$, and the right side decreases strictly with
$a$, from $p+1$ towards $p$. Let $a_p(\varepsilon)$ be the number of $a\ge0$ for which the inequality holds; it is $0$ as soon as $p^{1+\varepsilon}>p+1$.
\item The function $f$ attains its maximum. The number $N_\varepsilon=\prod_pp^{a_p(\varepsilon)}$ is a maximum point, and every maximum point divides it.
\end{enumerate}
\end{proposition}
\begin{proof}
(a) Since $\sigma(p^a)=\frac{p^{a+1}-1}{p-1}$, we have $\frac{f(p^{a+1})}{f(p^a)}=\frac{\sigma(p^{a+1})}{\sigma(p^a)\,p^{1+\varepsilon}}$, which gives the equivalence, and
$\frac{\sigma(p^{a+1})}{\sigma(p^a)}=\frac{p^{a+2}-1}{p^{a+1}-1}=p+\frac{p-1}{p^{a+1}-1}$, which decreases strictly from $p+1$ (at $a=0$) to $p$. So the inequality
holds for $a<a_p(\varepsilon)$ and fails for $a\ge a_p(\varepsilon)$: the sequence $a\mapsto f(p^a)$ increases weakly up to $a_p(\varepsilon)$, then decreases strictly. Its
largest maximum point is $a_p(\varepsilon)$, and every maximum point is at most $a_p(\varepsilon)$. If $p^{1+\varepsilon}>p+1$, the inequality already fails at $a=0$.

(b) Since $f$ is multiplicative, $f(n)=\prod_pf(p^{v_p(n)})$, a product of factors that can be chosen independently. For each $p$, $f(p^{v_p(n)})\le f(p^{a_p(\varepsilon)})$,
and for the finitely many $p$ with $a_p(\varepsilon)>0$ these are the only factors that can exceed $1$. So $f(n)\le f(N_\varepsilon)$ for every $n$, with equality if
and only if each $v_p(n)$ is a maximum point of $a\mapsto f(p^a)$, which forces $v_p(n)\le a_p(\varepsilon)$, that is $n\mid N_\varepsilon$.
\end{proof}

The numbers $N_\varepsilon$ are the \emph{colossally abundant} numbers. As $\varepsilon$ decreases, $a_p(\varepsilon)$ goes up by one exactly when $\varepsilon$ crosses
$\varepsilon_{p,a}=\log_p\frac{\sigma(p^{a+1})}{\sigma(p^a)}-1$, so $N_\varepsilon$ grows by a factor $p$ at each critical value. Sorting the critical values gives
\[
2,\ 6,\ 12,\ 60,\ 120,\ 360,\ 2520,\ 5040,\ 55\,440,\ 720\,720,\ 1\,441\,440,\ 4\,324\,320,\ 21\,621\,600,\ \dots
\]
They are the natural candidates for large values of $G$: $N_\varepsilon$ beats every $n$ in $\sigma(n)/n^{1+\varepsilon}$, so no smaller $n$ has a larger $\sigma(n)/n$,
and a larger $n$ needs a correspondingly larger $\sigma(n)/n$ to compete.

\section{Robin's criterion}

Compute $G(n)$ for small $n$ and you find values above $e^\gamma$: $G(3)\approx14.2$ because $\log\log3$ is tiny, and also $G(5040)\approx1.7910$. A
search shows that exactly $26$ integers $3\le n\le5040$ have $G(n)\ge e^\gamma$:
\[
3,4,5,6,8,9,10,12,16,18,20,24,30,36,48,60,72,84,120,180,240,360,720,840,2520,5040,
\]
and $5040$ is the last one. Figure~\ref{conjectures:fig:robin} shows $G$ along the colossally abundant numbers up to $10^{300}$: the $142$ of them beyond $5040$ all
satisfy $G(N)<e^\gamma$, and they approach $e^\gamma$ from below, very slowly (the last one is still $0.0092$ short).

\begin{figure}[ht]
\centering
\fig{conjectures-ch04-robin}
\caption{$G(n)=\frac{\sigma(n)}{n\log\log n}$ for the colossally abundant numbers $N\le10^{300}$ (blue dots) and for the integers $n\le5040$ with $G(n)\ge e^\gamma$
(red circles), on a doubly logarithmic scale in $n$; values above $2.3$ (the smallest $n$) are off the chart. The dashed line is $e^\gamma$. Beyond
$5040$ every colossally abundant number in range lies below it.}
\label{conjectures:fig:robin}
\end{figure}

\begin{remark}[Robin, 1984]\label{conjectures:rem:robin}
Guy Robin proved: the Riemann hypothesis holds if and only if $\sigma(n)<e^\gamma n\log\log n$ for every $n>5040$. Moreover, if the hypothesis is false,
the inequality fails for infinitely many $n$. One direction, under the hypothesis, refines Ramanujan's lost argument with explicit constants; the other
uses the zeros off the line to make $\prod_{p\le x}(1-\frac1p)^{-1}$ oscillate more than the hypothesis allows, as in Tome~5's proof of von Koch's
theorem. The proof is long and is not given here. Robin also observed that, if the hypothesis is false, a counterexample can be found among the
colossally abundant numbers; that is why the figure tests exactly those.
\end{remark}

In 2002 Jeffrey Lagarias turned Robin's criterion into a statement with no constants and no exceptions, using the harmonic numbers
$H_n=1+\frac12+\dots+\frac1n$.

\begin{lemma}\label{conjectures:lem:harmonic}
The sequence $H_n-\log n$ decreases strictly to Euler's constant $\gamma$. In particular $H_n>\log n+\gamma$ for every $n\ge1$.
\end{lemma}
\begin{proof}
The difference of consecutive terms is $\frac1{n+1}-\log\frac{n+1}n=\frac1{n+1}-\int_n^{n+1}\frac{\dd t}t<0$, because $\frac1t>\frac1{n+1}$ on $[n,n+1[$. The limit is $\gamma$
by definition (Tome~0, Chapter~5), and a strictly decreasing sequence stays above its limit.
\end{proof}

\begin{proposition}[Lagarias, the easy direction]\label{conjectures:prop:lagarias}
If $\sigma(n)<e^\gamma n\log\log n$ for every $n>5040$, then
\[
\sigma(n)\le H_n+e^{H_n}\log H_n\qquad\text{for every }n\ge1,
\]
with equality only for $n=1$. In particular, by Robin's theorem, the Riemann hypothesis implies Lagarias's inequality.
\end{proposition}
\begin{proof}
For $n>5040$, Lemma~\ref{conjectures:lem:harmonic} gives $e^{H_n}>e^\gamma n$ and $\log H_n>\log(\log n+\gamma)>\log\log n>0$. Hence
$H_n+e^{H_n}\log H_n>e^\gamma n\log\log n>\sigma(n)$. For $1\le n\le5040$ the inequality is a finite computation, carried out in the notebook: it holds
for every such $n$, with equality only at $n=1$, where both sides are $1$ (since $H_1=1$ and $\log1=0$).
\end{proof}

Lagarias also proved the converse: his inequality for all $n$ implies the Riemann hypothesis. It does not follow from Robin's criterion as stated,
because $H_n+e^{H_n}\log H_n$ is slightly larger than $e^\gamma n\log\log n$; it needs Robin's stronger statement that, if the hypothesis fails,
$\sigma(n)$ exceeds $e^\gamma n\log\log n$ by a definite margin infinitely often. So the Riemann hypothesis is equivalent to an inequality that a
first-year student can state and a computer can test for each $n$ in a fraction of a second.

\begin{remark}[other reformulations]
The Nyman--Beurling criterion (1950, 1955) says the hypothesis holds if and only if the constant function $1$ on $]0,1[$ is an $L^2$ limit of
combinations of the functions $x\mapsto\{\frac\theta x\}-\theta\{\frac1x\}$ with $0<\theta<1$. Li's criterion (1997) says it holds if and only if the numbers
$\lambda_n=\sum_\rho\bigl(1-(1-\frac1\rho)^n\bigr)$ are nonnegative for all $n\ge1$. Mertens' bound $M(x)=O(x^{1/2+\varepsilon})$ is another (Tome~5,
Exercise~7.3, one direction). Each turns the hypothesis into a statement in a different field, and none has been easier to prove.
\end{remark}

\section{Why it is hard}

There is no shortage of equivalent forms. What is missing is a reason. Three observations explain why the obvious routes fail.

\paragraph{The functional equation is not enough.} Harold Davenport and Hans Heilbronn constructed in 1936 a Dirichlet series, a combination of two
Dirichlet $L$-functions, that satisfies the same kind of functional equation as $\zeta$ and still has infinitely many zeros off the critical line, and even
to the right of $\operatorname{Re}s=1$. So any proof must use the Euler product, the multiplicativity that encodes unique factorisation. Analysis alone,
however sharp, does not see the primes.

\paragraph{The analogue is proved, but the proof does not transfer.} For curves over finite fields the analogue of the hypothesis is a theorem of
Andr\'e Weil (1948), proved in Tome~12, Chapter~13, by Bombieri's version of Stepanov's method. Every proof uses a geometric object, the curve, with a
Frobenius map acting on it: Weil's own proof uses a positivity argument in the intersection theory of the surface $C\times C$, Bombieri's an auxiliary
function that vanishes to high order at the points fixed by Frobenius. For $\mathbb Q$ nobody knows what the curve is, or what plays the role of
Frobenius.

\paragraph{The spectral dream.} Hilbert and P\'olya are said to have suggested, in conversation and without a written source, that the zeros are
$\frac12+i\lambda$ with $\lambda$ the eigenvalues of a self-adjoint operator, which would make them real. The statistics of the zeros, which match
those of random Hermitian matrices (Appendix~E), support the dream. Nobody has found the operator.

\section{Reading a paper: Lagarias and Conrey}

Jeffrey Lagarias, \emph{An elementary problem equivalent to the Riemann hypothesis} (American Mathematical Monthly, 2002), is ten pages long and can be
read in an afternoon. It states the inequality of Proposition~\ref{conjectures:prop:lagarias} as ``Problem E'', surveys the large values of $\sigma(n)$ and the colossally
abundant numbers (with a table up to $10^{18}$ and the open Alaoglu--Erd\H{o}s conjecture on their exponents), and then proves the equivalence from two
propositions of Robin and two elementary lemmas, the first of which is our easy direction. The idea worth stealing is the change of variable from
$\log n$ to $H_n$, which absorbs Euler's constant and the exception at $5040$ in one move.

Brian Conrey, \emph{The Riemann hypothesis} (Notices of the AMS, 2003), is a survey by one of the experts. Read its account of the zero counts and the
positive proportion results, its section on random matrices, and the list of approaches. Its tone, enthusiastic and sober at the same time, is a model
of how to write about an open problem.

\section{Open doors}

\begin{itemize}
\item Compute the first $10\,000$ Gram points and the proportion of $n$ for which Gram's law fails. Does it grow with $n$? Compare with the literature.
\item Lehmer found in 1956 two zeros near $t=7005$ that are unusually close together. Find the closest pairs of zeros among the first $10\,000$, normalised
by the mean spacing $\frac{2\pi}{\log(t/2\pi)}$.
\item Which $n\le10^7$ come closest to violating Lagarias's inequality? Are they colossally abundant?
\item The odd numbers: show by computer that $G(n)<e^\gamma$ for every odd $n$ with $9<n\le10^7$, and look up what is known for odd $n$ in general.
\end{itemize}

\section{Exercises}

\begin{exercise}[$\star$]\label{conjectures:ex:sigmaH}
Show that $\frac{\sigma(n)}n\le H_n\le1+\log n$ for $n\ge1$.
\end{exercise}

\begin{exercise}[$\star$]\label{conjectures:ex:robinfails}
Check that $\sigma(5040)=19\,344$, and that $\sigma(5040)>e^\gamma\cdot5040\cdot\log\log5040$. Which powers of $2$ violate Robin's inequality?
\end{exercise}

\begin{exercise}[$\star$]\label{conjectures:ex:lagsmall}
Check Lagarias's inequality by hand for $1\le n\le6$.
\end{exercise}

\begin{exercise}[$\star\star$]\label{conjectures:ex:landauphi}
Adapt the proof of Theorem~\ref{conjectures:thm:gronwall} to show Landau's theorem: $\displaystyle\liminf_{n\to\infty}\frac{\varphi(n)\log\log n}n=e^{-\gamma}$.
\end{exercise}

\begin{exercise}[$\star$]\label{conjectures:ex:cafirst}
Using Proposition~\ref{conjectures:prop:ca}, find $N_\varepsilon$ for $\varepsilon=\frac12$ and for $\varepsilon=\frac1{10}$.
\end{exercise}

\begin{exercise}[$\star\star$]\label{conjectures:ex:cacritical}
Compute the critical values $\varepsilon_{p,a}$ for $(p,a)=(2,0),(3,0),(2,1),(5,0),(2,2),(3,1),(7,0)$ and sort them. Deduce the first seven colossally abundant
numbers.
\end{exercise}

\begin{exercise}[$\star$]\label{conjectures:ex:Zgram}
Show that $\zeta(\frac12+ig_n)=(-1)^nZ(g_n)$, and that Gram's law at $g_n$ says $\operatorname{Re}\zeta(\frac12+ig_n)>0$.
\end{exercise}

\begin{exercise}[$\star\star$]\label{conjectures:ex:gramasym}
Using the formula for $\vartheta$, show that $g_n\sim\frac{2\pi n}{\log n}$, and that $N(g_n)=n+1+S(g_n)$. What does this say about the number of zeros between
consecutive Gram points, on average?
\end{exercise}

\begin{exercise}[$\star\star$]\label{conjectures:ex:gronwallsub}
Show that $G(n)\to0$ along the primes, and along the powers of $2$. Then show that for every $c\in[0,e^\gamma]$ there is a sequence $n_k\to\infty$ with
$G(n_k)\to c$.
\end{exercise}

\begin{exercise}[$\star$]\label{conjectures:ex:sigmamax}
Show that if $m\mid n$ then $\frac{\sigma(m)}m\le\frac{\sigma(n)}n$, with equality only for $m=n$. Deduce that every $n<N_\varepsilon$ with $\frac{\sigma(n)}n\ge\frac{\sigma(N_\varepsilon)}{N_\varepsilon}$ is
impossible: no smaller number is more abundant than a colossally abundant one.
\end{exercise}

\section{Hints}

\begin{hint}{conjectures:ex:sigmaH}
$\frac{\sigma(n)}n=\sum_{d\mid n}\frac1d$. For the second inequality compare with $\int_1^n\frac{\dd t}t$.
\end{hint}
\begin{hint}{conjectures:ex:robinfails}
$5040=2^4\cdot3^2\cdot5\cdot7$. For $2^k$, $\frac{\sigma(2^k)}{2^k}<2$.
\end{hint}
\begin{hint}{conjectures:ex:lagsmall}
$H_2=1.5$, $H_3\approx1.833$; $e^{H_n}\log H_n$ grows fast.
\end{hint}
\begin{hint}{conjectures:ex:landauphi}
$\frac n{\varphi(n)}=\prod_{p\mid n}(1-\frac1p)^{-1}$; use the primorials $\prod_{p\le x}p$ for the lower bound.
\end{hint}
\begin{hint}{conjectures:ex:cafirst}
List $\frac{\sigma(p^{a+1})}{\sigma(p^a)}$: for $p=2$ the values are $3,\frac73,\frac{15}7,\dots$.
\end{hint}
\begin{hint}{conjectures:ex:cacritical}
$\varepsilon_{2,0}=\log_23-1$.
\end{hint}
\begin{hint}{conjectures:ex:Zgram}
$e^{-i\vartheta(g_n)}=(-1)^n$.
\end{hint}
\begin{hint}{conjectures:ex:gramasym}
Invert $\vartheta(t)\approx\frac t2\log\frac t{2\pi}$, then use the formula for $N(T)$ at $T=g_n$.
\end{hint}
\begin{hint}{conjectures:ex:gronwallsub}
For the last part, multiply $n_x$ of the proof of Theorem~\ref{conjectures:thm:gronwall} by a large prime.
\end{hint}
\begin{hint}{conjectures:ex:sigmamax}
Use Proposition~\ref{conjectures:prop:ca}(b): $f(n)\le f(N_\varepsilon)$.
\end{hint}

\section{Solutions}

\begin{solution}{conjectures:ex:sigmaH}
\idea A sum over divisors is a sub-sum of the harmonic series.
\writeup $\frac{\sigma(n)}n=\sum_{d\mid n}\frac nd\cdot\frac1n=\sum_{d\mid n}\frac1d$, because $d\mapsto\frac nd$ permutes the divisors. The divisors are among
$1,\dots,n$, so the sum is at most $H_n$. Since $\frac1k\le\int_{k-1}^k\frac{\dd t}t$ for $k\ge2$, $H_n\le1+\int_1^n\frac{\dd t}t=1+\log n$.
\end{solution}

\begin{solution}{conjectures:ex:robinfails}
\idea Multiplicativity of $\sigma$.
\writeup $\sigma(5040)=\sigma(16)\sigma(9)\sigma(5)\sigma(7)=31\cdot13\cdot6\cdot8=19\,344$. And $\log5040=8.5252$, $\log\log5040=2.1430$, so
$e^\gamma\cdot5040\cdot2.1430=19\,237.1<19\,344$. For $n=2^k$, $G(2^k)=\frac{2-2^{-k}}{\log(k\log2)}$, which is at least $e^\gamma$ exactly for $k=2,3,4$ (the values are
$5.36$, $2.56$, $1.90$) and below for $k\ge5$ ($1.58$ at $k=5$, decreasing after). So $4,8,16$ violate the inequality, as in the list of the chapter.
\end{solution}

\begin{solution}{conjectures:ex:lagsmall}
\idea Plug in.
\writeup With $L_n=H_n+e^{H_n}\log H_n$: $L_1=1=\sigma(1)$; $L_2=1.5+e^{1.5}\log1.5=1.5+4.482\cdot0.405=3.32>3$; $L_3=1.833+6.255\cdot0.606=5.62>4$;
$L_4=2.083+8.03\cdot0.734=7.98>7$; $L_5=2.283+9.81\cdot0.826=10.38>6$; $L_6=2.45+11.59\cdot0.896=12.84>12$. The inequality holds, with equality only at $n=1$; $n=6$
comes closest in this range.
\end{solution}

\begin{solution}{conjectures:ex:landauphi}
\idea The same two halves as Gronwall's theorem, with $\frac n{\varphi(n)}$ in place of $\frac{\sigma(n)}n$.
\writeup Since $\frac n{\varphi(n)}=\prod_{p\mid n}(1-\frac1p)^{-1}$, the upper-bound half of the proof of Theorem~\ref{conjectures:thm:gronwall}, which in fact bounded this product, gives
$\frac n{\varphi(n)}\le(1+\varepsilon)e^\gamma\log\log n$ for large $n$, that is $\frac{\varphi(n)\log\log n}n\ge\frac{e^{-\gamma}}{1+\varepsilon}$. For the other direction take the
primorials $m_x=\prod_{p\le x}p$: $\frac{m_x}{\varphi(m_x)}=M(x)\sim e^\gamma\log x$ and $\log\log m_x=\log\theta(x)=\log x+o(1)$, so
$\frac{\varphi(m_x)\log\log m_x}{m_x}\to e^{-\gamma}$.
\end{solution}

\begin{solution}{conjectures:ex:cafirst}
\idea Count, for each prime, how many ratios $\frac{\sigma(p^{a+1})}{\sigma(p^a)}$ are at least $p^{1+\varepsilon}$.
\writeup For $\varepsilon=\frac12$: $2^{1.5}=2.83\le3$ but $>\frac73$, so $a_2=1$; $3^{1.5}=5.20>4$, so $a_3=0$, and larger primes give $0$ too. So $N_{1/2}=2$.
For $\varepsilon=\frac1{10}$: $2^{1.1}=2.1435$, and the ratios $3$, $\frac73=2.333$ are larger while $\frac{15}7=2.1429$ is smaller, so $a_2=2$; $3^{1.1}=3.348$ is below
$4$ but above $\frac{13}4$, so $a_3=1$; $5^{1.1}=5.873\le6$ but $>\frac{31}6$, so $a_5=1$; $7^{1.1}=8.50>8$, so $a_7=0$, and $p^{1.1}>p+1$ for all larger $p$. So
$N_{1/10}=4\cdot3\cdot5=60$.
\end{solution}

\begin{solution}{conjectures:ex:cacritical}
\idea Each critical value adds one prime factor.
\writeup $\varepsilon_{p,a}=\frac{\log(\sigma(p^{a+1})/\sigma(p^a))}{\log p}-1$: $\varepsilon_{2,0}=\log_23-1=0.585$; $\varepsilon_{3,0}=\log_34-1=0.262$; $\varepsilon_{2,1}=\log_2\frac73-1=0.222$;
$\varepsilon_{5,0}=\log_56-1=0.113$; $\varepsilon_{2,2}=\log_2\frac{15}7-1=0.0995$; $\varepsilon_{3,1}=\log_3\frac{13}4-1=0.0729$; $\varepsilon_{7,0}=\log_78-1=0.0686$. One also needs
$\varepsilon_{2,3}=\log_2\frac{31}{15}-1=0.0473$ and $\varepsilon_{11,0}=\log_{11}12-1=0.0363$, which are smaller. In decreasing order the factors are $2,3,2,5,2,3,7$, and the
partial products are $2,6,12,60,120,360,2520$.
\end{solution}

\begin{solution}{conjectures:ex:Zgram}
\idea At a Gram point the phase factor is a sign.
\writeup $\zeta(\frac12+it)=e^{-i\vartheta(t)}Z(t)$, and $e^{-i\vartheta(g_n)}=e^{-in\pi}=(-1)^n$. So $\zeta(\frac12+ig_n)=(-1)^nZ(g_n)$, a real number, and $(-1)^nZ(g_n)>0$ means
that this real number, which is its own real part, is positive.
\end{solution}

\begin{solution}{conjectures:ex:gramasym}
\idea $\vartheta$ is a smooth, slowly accelerating clock; the Gram points are its ticks.
\writeup $\vartheta(g_n)=n\pi$ and $\vartheta(t)=\frac t2\log\frac t{2\pi}+O(t)$, more precisely $\frac t2\log t\,(1+O(\frac1{\log t}))$. So $n\pi\sim\frac{g_n}2\log g_n$, hence
$\log g_n\sim\log n$ and $g_n\sim\frac{2\pi n}{\log n}$. By Theorem~6.9 of Tome~5, $N(g_n)=\frac{\vartheta(g_n)}\pi+1+S(g_n)=n+1+S(g_n)$ (when $g_n$ is not the ordinate of a
zero). So the number of zeros up to $g_n$ differs from $n+1$ by $S(g_n)=O(\log g_n)$: on average there is exactly one zero per Gram interval, and the
irregularities are measured by $S$. Since $S(t)$ has mean $0$ (Littlewood), Gram intervals with two zeros and with none balance out.
\end{solution}

\begin{solution}{conjectures:ex:gronwallsub}
\idea $\frac{\sigma(n)}n$ is bounded on these sequences while $\log\log n\to\infty$; multiplying by a big prime lowers $G$ continuously enough.
\writeup For a prime $p$, $\frac{\sigma(p)}p=1+\frac1p\le\frac32$; for $2^k$, $\frac{\sigma(2^k)}{2^k}<2$. In both cases $\frac{\sigma(n)}n$ is bounded and $G\to0$.
For $c=e^\gamma$ use the numbers $n_x$ of the proof. For $0<c<e^\gamma$, fix a large $x$, let $L=\log n_x$ and consider $G(n_xq)$ for the primes $q>x$:
\[
G(n_xq)=\frac{\sigma(n_x)}{n_x}\Bigl(1+\frac1q\Bigr)\frac1{\log(L+\log q)}.
\]
For the first prime $q>x$ (which is at most $2x$, by Bertrand) this is $G(n_x)(1+o(1))$, close to $e^\gamma$; as $q\to\infty$ it tends to $0$. Between consecutive
primes $q<q'\le2q$, the factor $\log(L+\log q)$ changes by at most $\log\bigl(1+\frac{\log2}L\bigr)\le\frac{\log2}L$ in absolute value and $1+\frac1q$ changes by less than $\frac1x$,
so $G$ changes by a factor $1+O(\frac1L+\frac1x)$. Hence some $q$ gives $|G(n_xq)-c|=O(\frac1L+\frac1x)+o(1)$, and letting $x\to\infty$ gives a sequence with
limit $c$.
\end{solution}

\begin{solution}{conjectures:ex:sigmamax}
\idea Divisors of $m$ are divisors of $n$.
\writeup $\frac{\sigma(m)}m=\sum_{d\mid m}\frac1d\le\sum_{d\mid n}\frac1d=\frac{\sigma(n)}n$, and the inequality is strict if $m<n$, because $n$ itself is a divisor of $n$ but not of $m$. For
the second part, if $n<N_\varepsilon$ and $\frac{\sigma(n)}n\ge\frac{\sigma(N_\varepsilon)}{N_\varepsilon}$, then $f(n)=\frac{\sigma(n)}n\cdot\frac1{n^\varepsilon}>\frac{\sigma(N_\varepsilon)}{N_\varepsilon}\cdot\frac1{N_\varepsilon^\varepsilon}=f(N_\varepsilon)$,
contradicting Proposition~\ref{conjectures:prop:ca}(b).
\end{solution}

\section*{Sources and further reading}

\begin{itemize}
\item B.~Riemann, \emph{Ueber die Anzahl der Primzahlen unter einer gegebenen Gr\"osse}, Monatsberichte der K\"oniglich Preussischen Akademie der
Wissenschaften zu Berlin, November 1859, 671--680 (the epigraph: ``es ist sehr wahrscheinlich, dass alle Wurzeln reell sind. Hiervon w\"are allerdings
ein strenger Beweis zu w\"unschen'', with the rest of the sentence in our words). C.~L.~Siegel, \emph{\"Uber Riemanns Nachla\ss{} zur analytischen
Zahlentheorie}, Quellen und Studien zur Geschichte der Mathematik, Astronomie und Physik, Abt.~B, 2 (1932), 45--80.
\item J.-P.~Gram, \emph{Note sur les z\'eros de la fonction $\zeta(s)$ de Riemann}, Acta Mathematica 27 (1903), 289--304. J.~I.~Hutchinson, \emph{On the
roots of the Riemann zeta function}, Transactions of the AMS 27 (1925), 49--60. A.~M.~Turing, \emph{Some calculations of the Riemann zeta-function},
Proceedings of the London Mathematical Society (3) 3 (1953), 99--117.
\item P.~Moree, I.~Petrykiewicz and A.~Sedunova, \emph{A computational history of prime numbers and Riemann zeros}, arXiv:1810.05244 (2018), for the
history of the zero counts. D.~Platt and T.~Trudgian, \emph{The Riemann hypothesis is true up to $3\cdot10^{12}$}, Bulletin of the London Mathematical
Society 53 (2021).
\item T.~H.~Gronwall, \emph{Some asymptotic expressions in the theory of numbers}, Transactions of the AMS 14 (1913), 113--122.
L.~Alaoglu and P.~Erd\H{o}s, \emph{On highly composite and similar numbers}, Transactions of the AMS 56 (1944), 448--469.
\item S.~Ramanujan, \emph{Highly composite numbers}, annotated by J.-L.~Nicolas and G.~Robin, The Ramanujan Journal 1 (1997), 119--153. J.-L.~Nicolas and
J.~Sondow, \emph{Ramanujan, Robin, highly composite numbers, and the Riemann hypothesis}, arXiv:1211.6944, for the story of the lost pages.
\item G.~Robin, \emph{Grandes valeurs de la fonction somme des diviseurs et hypoth\`ese de Riemann}, Journal de Math\'ematiques Pures et Appliqu\'ees 63
(1984), 187--213.
\item J.~C.~Lagarias, \emph{An elementary problem equivalent to the Riemann hypothesis}, American Mathematical Monthly 109 (2002), 534--543.
\item J.~B.~Conrey, \emph{The Riemann hypothesis}, Notices of the AMS 50 (2003), 341--353.
\item X.-J.~Li, \emph{The positivity of a sequence of numbers and the Riemann hypothesis}, Journal of Number Theory 65 (1997), 325--333.
A.~Beurling, \emph{A closure problem related to the Riemann zeta-function}, Proceedings of the National Academy of Sciences 41 (1955), 312--314.
\item T.~Trudgian, \emph{On the success and failure of Gram's law and the Rosser rule}, Acta Arithmetica 148 (2011).
\item H.~Davenport and H.~Heilbronn, \emph{On the zeros of certain Dirichlet series}, Journal of the London Mathematical Society 11 (1936).
\end{itemize}
